\documentclass[12pt, a4paper]{ctexart}

% ---- 基础宏包 ----
\usepackage{amsmath, amssymb, amsthm}
\usepackage{graphicx}
\usepackage{booktabs}
\usepackage{longtable}
\usepackage{array}
\usepackage{multirow}
\usepackage{xcolor}
\usepackage{hyperref}
\usepackage{geometry}
\usepackage{setspace}
\usepackage{caption}
\usepackage{subcaption}
\usepackage{enumitem}
\usepackage{float}

% ---- 页面设置 ----
\geometry{
  a4paper,
  top=2.54cm, bottom=2.54cm,
  left=3.17cm, right=3.17cm
}

\onehalfspacing

% ---- 字体 ----
\setCJKmainfont{SimSun}        % 宋体（或改为 Source Han Serif CN 等）
\setCJKsansfont{SimHei}        % 黑体
\setmainfont{Times New Roman}

% ---- 章节编号深度 ----
\setcounter{secnumdepth}{3}
\setcounter{tocdepth}{3}

\hypersetup{
  colorlinks=true,
  linkcolor=black,
  citecolor=blue,
  urlcolor=blue
}

\begin{document}



\section{第四章 面向多场景验证的多尺度物理一致性锂离子电池SOH估计与异常识别方法}


\section{4.1 引言}

锂离子电池健康状态（State of Health, SOH）的准确评估是电池管理系统（Battery Management System, BMS）实现安全运行、寿命管理与故障预警的重要基础。在长期服役过程中，电池会经历活性物质损失、电解液分解、SEI膜增厚和极化加剧等复杂退化过程，其宏观表现为可用容量衰减、内阻增长及充放电响应特征变化。由于电池类型、制造一致性、运行工况和环境条件存在差异，SOH估计不仅需要较高的预测精度，还需要在跨电池、跨阶段和标定不完整条件下保持稳定的泛化能力。


现有SOH估计方法主要包括实验测量法、基于模型的方法和数据驱动方法。安时积分法、容量测试和等效电路模型等方法具有较强的物理可解释性，但通常依赖严格测试条件、完整容量标定或复杂参数辨识过程，难以适应长期在线运行中的非线性退化行为。相比之下，数据驱动方法能够从电压、电流、温度及充电片段等可测信号中学习状态特征与老化程度之间的映射关系，因而逐渐成为当前研究的重要方向。


国外早期数据驱动研究多采用“健康因子提取—机器学习回归”的技术路线，即从容量、内阻、增量容量曲线和恒流/恒压充电片段中提取特征，并结合支持向量回归、高斯过程回归、随机森林等模型进行SOH或寿命预测。Severson等利用早期循环数据实现了电池寿命预测，证明了机器学习方法在电池退化建模中的潜力[1]；Roman等从机器学习角度总结了SOC、SOH和RUL联合估计的研究进展[2]。随着深度学习的发展，CNN、LSTM、CNN-BiLSTM、Transformer及迁移学习方法被广泛用于SOH/RUL估计，能够分别刻画充电曲线局部模式、跨循环时序依赖以及跨电池差异[3-11]。这些研究在公开数据集上取得了较高精度，但其性能通常依赖较完整的标签信息和相对稳定的数据分布。


与国外研究趋势相一致，国内学者也围绕锂离子电池SOH估计开展了大量研究，主要集中在健康因子提取、深度时序网络、融合模型和工程在线估计等方向。一类工作从CC-CV充电过程、电压平台、容量增量曲线及温度统计量中构造健康特征，并结合BP神经网络、支持向量机、随机森林或集成学习模型估计SOH；另一类工作则采用LSTM、GRU、CNN-LSTM、1DCNN-LSTM及注意力机制等结构，以提高模型对非线性退化过程的表征能力[12-15]。总体来看，国内研究较重视可在线提取特征与工程实现便利性，为实际BMS应用提供了较多参考，但多数研究仍以完整实验室循环数据或随机划分样本作为主要验证方式，对跨电池泛化、标签稀疏和无标签区间训练机制的讨论仍相对不足。


从验证范式看，公开数据集和跨电池验证正在成为SOH估计研究的重要基础。Battery Data Genome强调电池数据共享、元数据规范和跨机构可复现实验的重要性[16]；HUST公开数据集提供了77节LFP/石墨电池在统一充电协议和多阶段放电协议下的全生命周期数据，为研究个体差异、跨电池泛化和工况变化提供了条件[17]。然而，公开数据集并不等同于真实BMS部署环境。许多研究在建模时直接使用完整循环SOH作为监督信号，而实际系统中往往只能连续记录运行特征，SOH标签则依赖间歇标定获得。因此，仅在完整标签条件下追求较低误差，尚不足以说明模型具备低标定成本下的工程可用性。后续研究有必要面向标签稀疏、跨电池差异和工况变化等更接近实际应用的条件，构建兼顾预测精度、泛化能力和工程可部署性的SOH估计方法。


为提升深度模型的物理可信度，物理信息机器学习和物理约束神经网络逐渐被引入电池退化建模。相关研究通过将退化单调性、容量边界、动力学残差或增量容量特征等先验嵌入学习过程，提高了SOH预测的可解释性和外推稳定性[18-21]。这些工作表明，将电化学退化规律融入深度学习模型，有助于缓解纯数据驱动方法在小样本和分布外场景下的不稳定性。然而，现有物理信息SOH研究大多仍默认SOH标签可用于连续训练，较少将“运行特征可观测但SOH标签缺失”的部分生命周期监督结构作为核心问题进行建模。


综合来看，现有SOH估计研究仍存在以下不足。首先，部分方法采用样本级随机划分或完整生命周期标签训练，训练集与测试集之间可能存在同一电池不同循环的强相关性，从而高估模型泛化能力。其次，深度网络多关注拟合精度提升，对SOH标签稀疏和监督区间不完整条件下模型如何获得有效训练信号讨论不足。再次，物理约束常被作为普通正则项使用，其在无标签循环中的弱监督作用尚未得到充分解释。最后，注意力、伪标签、不确定性估计和自适应权重等复杂模块虽然能够增强模型表达能力，但在低标签条件下也可能带来校准不足、权重漂移或局部过拟合等问题，其工程收益边界仍需进一步分析。


基于上述问题，本文提出一种面向部分生命周期监督的多尺度物理一致性CNN-LSTM模型（MS-PI-CNNLSTM）。该模型以多尺度卷积提取不同时间尺度下的退化特征，以LSTM建模跨循环状态演化，并通过软单调物理约束限制SOH预测轨迹中的非物理回弹。在缺少SOH标签的循环中，数据误差项无法提供直接监督，而退化方向约束仍可对模型输出形成结构性引导。因此，本文将物理一致性约束视为部分生命周期监督场景下的弱监督信号，而不仅是提高预测精度的普通正则化项。


相较于已有研究，本文的特点主要体现在三个方面：第一，在问题设定上，显式构建部分生命周期监督框架，保留完整运行特征轨迹，仅屏蔽缺失SOH标签的数据误差项，以贴近实际BMS中运行数据连续可采集而容量标定稀疏的应用场景；第二，在模型结构上，采用多尺度CNN-LSTM增强有限标签条件下的退化表征能力，避免单一卷积感受野难以兼顾局部波动与长期趋势的问题；第三，在训练机制上，将软单调约束用于无标签循环的结构性引导，明确物理一致性约束在部分监督SOH估计中的作用。进一步地，本文在模型评价中不仅关注MA和RMSE，也关注预测轨迹的趋势一致性、跨电池稳定性以及模型结构的可解释性和可复现性，从而更全面地评估其在工程场景中的适用性。



\section{4.2 数据准备与特征工程}


\subsection{4.2.1 原始HUST数据集与特征空间构建}

本研究基于华中科技大学（HUST）公开锂离子电池全生命周期退化数据集，选取其中77组LFP/石墨电池单元作为研究对象。每个电池单元的标称容量为1.1 Ah，标称电压为3.3 V，所有样本在30 °C恒温环境下运行，并采用统一的CC-CV充电协议。放电阶段施加不同的多阶段放电工况，用于模拟实际应用中负载变化对退化轨迹的影响[17,22]。


相较于已有研究，本文的优势主要体现在三个方面。第一，在验证路线方面，本文不再直接将算法置于不完整标定场景中讨论，而是采用由易到难的三阶段验证逻辑：首先在HUST公开正常退化数据集上验证模型的基础SOH估计能力；其次在自建不完整标定数据集上考察模型在低标定成本条件下的适用性；最后在自建异常电池数据集上进一步评估SOH预测轨迹对异常退化和故障风险的识别能力。第二，在模型结构上，本文采用多尺度CNN-LSTM提取短程、中程和长程退化特征，避免单一卷积感受野难以兼顾局部波动和长期趋势的问题。第三，在训练机制上，本文将软单调约束解释为退化轨迹的物理一致性约束；在完整标定数据上，它用于提高预测轨迹的物理可信度，在不完整标定数据上，它进一步为缺失标定区间提供结构性弱监督信号。


% 图4-1用于展示所有电池在全寿命循环过程中的容量随循环次数变化情况。不同电池在寿命长度、衰减速率和局部波动幅度上存在明显差异，说明该数据集能够覆盖较为丰富的退化形态，为跨电池SOH估计和泛化能力验证提供了数据基础。

本章的研究重点不是简单堆叠多个深度学习模块，而是围绕“正常退化估计—不完整标定估计—异常退化识别”的工程递进关系构建一条可解释的建模路线。首先，利用HUST公开数据集建立标准SOH估计基准，检验MS-PI-CNNLSTM在正常退化电池上的精度和泛化能力；其次，将模型迁移到自建不完整标定数据集，分析SOH标签缺失时多尺度表征和物理一致性约束的作用；最后，基于自建异常电池数据集或异常注入数据，考察SOH预测轨迹是否能够为故障识别提供额外信息。这样安排能够避免论文一开始就落入不完整监督这一较复杂设定，使读者先理解模型在正常退化条件下的有效性，再逐步接受低标定成本和异常识别场景中的扩展价值。


% 表4-1  CC-CV充电循环中提取的统计特征与动力学特征定义（具体特征名称可按最终代码口径回填）


\begin{table}[H]
  \centering
  \begin{tabular}{lll}
  \toprule
  类别 & 特征示例 & 物理含义 \\
  \midrule
  电压统计特征 & 均值、标准差、偏度、峰度 & 反映充电电压曲线形态及极化变化 \\
  电流统计特征 & 均值、标准差、偏度、峰度 & 反映CV阶段电流衰减特性 \\
  阶段动力学特征 & CC时长、CV时长、充入电量 & 反映充电过程持续时间与可接受电量变化 \\
  容量相关特征 & 归一化容量或派生健康因子 & 刻画SOH退化趋势 \\
  \bottomrule
  \end{tabular}
  \caption{}
  \label{tab:}
\end{table}

需要说明的是，部分生命周期监督在本文中是第二阶段验证场景，而不是唯一研究起点。完整或较完整标定的正常退化数据用于回答“模型是否具备可靠的SOH估计能力”；不完整标定数据用于回答“在标定成本受限时模型能否维持可用性能”；异常电池数据用于回答“SOH预测轨迹能否为故障识别提供早期提示”。这三个问题分别对应电池健康管理中的基础估计、低成本标定和安全预警三个层面，构成本文实验组织的主线。



\begin{figure}[H]
  \centering
  % \includegraphics[width=0.8\textwidth]{figures/placeholder}
  \caption{图4-2  多源特征与容量退化之间的相关性矩阵（图片占位）}
  \label{fig:}
\end{figure}

4.2.1 HUST公开正常退化数据集与特征空间构建


需要说明的是，相关性分析只能揭示单个特征与容量之间的线性关联，无法完整刻画充电过程中的多变量耦合关系。例如，某些电压统计量与容量之间的单变量相关性可能并不突出，但其与CV阶段电流衰减特征组合后，可能反映极化增强或内部阻抗变化。深度模型的优势在于能够学习这种非线性组合关系。因此，本文在特征构建阶段不过度依赖单一相关性阈值，而是通过后续多尺度卷积和时序建模自动提取有效退化模式[27-28]。


HUST数据集中的电池均采用统一充电协议，但放电工况存在差异，这使得充电特征既包含由电池本体退化引起的长期变化，也可能受到前一放电阶段工况差异的间接影响。对于SOH估计模型而言，这种差异既是挑战，也是检验泛化能力的重要条件。若模型只能记忆某些固定工况下的特征模式，则在未见电池或不同退化阶段上容易失效；若模型能够从多粒度特征中捕捉稳定退化规律，则更有可能适应复杂服役条件[22]。


从硕士论文写作角度，本节还需要明确数据集选择与研究问题之间的对应关系。HUST数据集的优势不只在于电池数量较多，更在于其包含统一充电协议和差异化放电协议的组合：统一充电协议保证特征提取口径相对一致，差异化放电协议则引入个体退化差异和工况扰动。Battery Data Genome相关工作强调，电池数据研究需要同时关注数据规模、元数据规范、实验协议和跨数据源可复现性[16]。因此，本文选择HUST数据集并采用电池级划分，可以在相对受控的实验条件下讨论跨电池泛化，同时为后续转化到真实BMS运行数据提供较清晰的验证起点。



\subsection{4.2.2 基于部分生命周期监督的退化样本构建}

在实际BMS应用中，SOH标定通常依赖周期性容量测试、离线诊断或维护窗口，难以在电池完整生命周期内连续获得。因此，本文不假设每个充电循环均具有SOH标签，而是将SOH估计建模为部分生命周期监督任务。在该任务中，训练电池的完整运行特征序列可用，但仅有部分连续生命周期区间具有SOH监督标签。近年来半监督SOH估计、迁移学习和无标签充电数据利用研究均指出，真实应用中无标签运行数据远多于容量标签，如何有效利用这些未标注数据已成为电池健康估计的重要方向[29-31]。


公式4-1  X\textbackslash{}textasciicircum\{\}(i)=\textbackslash{}\{x\textbackslash{}\_1\textbackslash{}textasciicircum\{\}(i),...,x\textbackslash{}\_Ti\textbackslash{}textasciicircum\{\}(i)\textbackslash{}\},  y\textbackslash{}textasciicircum\{\}(i)=\textbackslash{}\{y\textbackslash{}\_1\textbackslash{}textasciicircum\{\}(i),...,y\textbackslash{}\_Ti\textbackslash{}textasciicircum\{\}(i)\textbackslash{}\}


设第 i 个电池的完整生命周期长度为 Ti，其充电循环特征序列和SOH序列分别如公式4-1所示。训练阶段仅选取其中一个或多个连续子区间作为可用监督标签，其余时间步对应的SOH不参与数据误差项计算。为描述标签可用性，定义监督掩码如公式4-2所示。


公式4-2  m\textbackslash{}\_t\textbackslash{}textasciicircum\{\}(i)=1, if t in Omega\textbackslash{}\_i;  m\textbackslash{}\_t\textbackslash{}textasciicircum\{\}(i)=0, if t not in Omega\textbackslash{}\_i


其中，m\textbackslash{}\_t\textbackslash{}textasciicircum\{\}(i)=1表示第i个电池第t个循环具有SOH标签，m\textbackslash{}\_t\textbackslash{}textasciicircum\{\}(i)=0表示该循环仅有运行特征而无SOH监督。本文采用Label Masking而非Sample Dropping：无标签循环不会被删除，其运行特征仍参与模型前向传播，只是不计算数据监督损失。该设定保留了完整生命周期特征轨迹，更符合实际BMS中运行数据连续可得、健康标定稀疏可得的工程条件。



\begin{figure}[H]
  \centering
  % \includegraphics[width=0.8\textwidth]{figures/placeholder}
  \caption{图4-3  部分生命周期监督条件下训练电池的标签覆盖示意图（图片占位）}
  \label{fig:}
\end{figure}

从硕士论文写作角度，本节还需要明确数据集选择与研究问题之间的对应关系。HUST数据集的优势不只在于电池数量较多，更在于其包含统一充电协议和差异化放电协议的组合：统一充电协议保证特征提取口径相对一致，差异化放电协议则引入个体退化差异和工况扰动。Battery Data Genome相关工作强调，电池数据研究需要同时关注数据规模、元数据规范、实验协议和跨数据源可复现性[16]。因此，本文首先选择HUST公开数据集作为正常退化基准，用于验证模型在较规范公开数据上的基础估计能力；随后再引入自建数据集讨论不完整标定和异常识别问题。


4.2.2 自建不完整标定与异常电池数据集构建


在完成HUST公开正常退化数据集上的基础验证后，本文进一步构建自建数据集，用于模拟和评估更接近工程应用的两类场景。第一类为不完整标定数据集，即电池运行特征能够连续采集，但SOH标定仅在部分维护窗口或测试区间内获得；第二类为异常电池数据集，即电池在服役过程中出现突发老化、退化拐点、单体不均衡、析锂或内阻异常增长等非典型退化行为。前者用于验证低标定成本下的SOH估计能力，后者用于验证SOH预测轨迹在故障识别中的辅助价值。


本文的Label Masking与传统半监督伪标签策略存在区别。半监督方法通常需要对无标签样本生成伪标签，或通过协同训练、迁移组件分析等方式扩大有效训练集[29-31]；而本文首先不假设伪标签一定可靠，而是让无标签循环保留在前向传播和物理约束计算中，通过退化方向、边界范围和序列平滑等低风险先验参与训练。这种处理更保守，也更适合当前实验尚未全面验证伪标签校准质量的阶段。换言之，本文把“无标签样本如何参与训练”分解为两个层次：数值标签不足时不强行制造标签，趋势规律明确时利用物理一致性提供弱监督。



\subsection{4.2.3 数据预处理与电池级验证划分}

其中，m\textbackslash{}\_t\textbackslash{}textasciicircum\{\}(i)=1表示第i个电池第t个循环具有SOH标定值，m\textbackslash{}\_t\textbackslash{}textasciicircum\{\}(i)=0表示该循环仅有运行特征而缺少同步SOH标定。对于自建不完整标定数据集，本文采用Label Masking而非Sample Dropping：无标定循环不会被删除，其运行特征仍参与模型前向传播，只是不计算数据监督损失。该设定保留了完整生命周期特征轨迹，更符合实际BMS中运行数据连续可得、健康标定稀疏可得的工程条件。


在数据划分策略上，本文采用严格的电池级别划分，而非样本级随机划分。训练集、验证集和测试集按电池个体划分，确保测试电池在训练阶段完全不可见，从而更真实地评估模型对未见电池个体的泛化能力。



\begin{figure}[H]
  \centering
  % \includegraphics[width=0.8\textwidth]{figures/placeholder}
  \caption{图4-3  自建不完整标定数据集中的标签覆盖示意图（图片占位）}
  \label{fig:}
\end{figure}

% 图4-3用于展示自建不完整标定数据集中不同电池在生命周期维度上的SOH标签覆盖区间。不同电池的标定区间可覆盖早期、中期或后期退化阶段，从群体层面提供跨阶段信息，同时避免依赖单一电池完整退化轨迹。预测阶段模型仍对任意给定循环输出对应SOH，不完整标定仅改变训练阶段的标签可见性，不改变部署时的推理方式。

这种设定的难点在于，模型不能简单依赖某一电池的完整退化轨迹进行外推，而需要从多个电池的局部标定片段中归纳出跨个体共享的退化映射关系。例如，某些电池可能只在早期健康阶段具有标签，另一些电池可能只在中后期退化阶段具有标签，模型必须在群体层面整合这些片段信息。若没有物理一致性约束，无标签区间的预测值只受到相邻有标签片段和模型参数共享的间接影响，容易出现局部漂移或不合理波动。



\begin{table}[H]
  \centering
  \begin{tabular}{lll}
  \toprule
  项目 & 设置 & 说明 \\
  \midrule
  划分对象 & 电池个体 & 避免同一电池循环同时出现在训练集和测试集 \\
  划分比例 & 60\% / 20\% / 20\% & 分别用于训练、验证和测试 \\
  监督方式 & Label Masking & 仅屏蔽标签损失，不删除无标签样本 \\
  测试目标 & 未见电池SOH估计 & 评估跨电池泛化能力 \\
  \bottomrule
  \end{tabular}
  \caption{}
  \label{tab:}
\end{table}

本文的Label Masking与传统半监督伪标签策略存在区别。半监督方法通常需要对无标签样本生成伪标签，或通过协同训练、迁移组件分析等方式扩大有效训练集[29-31]；而本文首先不假设伪标签一定可靠，而是让无标签循环保留在前向传播和物理约束计算中，通过退化方向、边界范围和序列平滑等低风险先验参与训练。这种处理更保守，也更适合当前实验尚未全面验证伪标签校准质量的阶段。



\section{异常电池数据集则用于第三阶段验证。与不完整标定数据不同，异常数据关注的不是SOH标签是否稀疏，而是电池退化轨迹是否偏离正常模式。本文将异常识别定位为SOH估计模型的工程扩展：在不重新训练专用故障诊断模型的前提下，利用已训练SOH模型输出的预测轨迹、局部趋势偏差和输入特征偏离程度，判断电池是否出现异常退化风险。相关异常类型和检测信号将在4.4.7节进一步说明。}

4.2.3 数据预处理与跨场景验证划分



\subsection{4.3.1 整体框架概述}

该框架包括三个组成部分。第一，多尺度一维卷积特征提取模块通过不同感受野并行提取多粒度充电特征；第二，LSTM时序建模层对跨循环退化状态进行编码；第三，物理一致性弱监督模块在损失函数层面嵌入退化单调性先验，使未标注循环也能对模型训练产生约束作用。



\begin{figure}[H]
  \centering
  % \includegraphics[width=0.8\textwidth]{figures/placeholder}
  \caption{图4-4  MS-PI-CNNLSTM整体框架示意图（图片占位）}
  \label{fig:}
\end{figure}

此外，不完整标定任务要求在数据划分之后再生成监督掩码，而不能先在全数据上随机抽取有标签样本后再划分训练测试集。原因在于，如果标签掩码生成过程跨越训练集和测试集，可能会使测试电池的生命周期分布信息间接影响训练设置，降低验证的严格性。本文建议固定电池级划分后，在训练集内部根据监督比例和覆盖区间生成Label Masking矩阵，验证集用于选择超参数和早停，测试集只用于最终泛化评价。



\subsection{4.3.2 多尺度卷积特征提取模块}

锂离子电池退化信号同时包含不同时间尺度的信息。短程尺度上，相邻循环的局部波动反映即时电化学响应变化；中程尺度上，连续若干循环的特征变化体现阶段性极化和活性物质损失；长程尺度上，整体容量衰减趋势揭示宏观退化路径。固定卷积核的单尺度CNN难以同时捕获这些信息。多尺度卷积和时间卷积网络在通用时序建模中已经被证明能够通过不同感受野提取局部与长程模式，InceptionTime和TCN等结构也说明了并行卷积或扩张卷积对于多时间尺度序列特征学习的价值[11,32]。


为此，本文设计三路并行多尺度一维卷积模块，分别采用kernel=3、kernel=7和kernel=15的卷积分支。短程分支聚焦相邻循环的局部动态，中程分支刻画阶段性变化，长程分支捕获更宽时间范围内的退化趋势。各分支输出在通道维度拼接后，通过1×1卷积融合至统一特征空间，再输入LSTM层。与单纯加深网络不同，并行多尺度结构能够显式保留不同退化粒度的特征通道，便于后续LSTM在循环维度上组织状态演化。


MS-PI-CNNLSTM的设计服务于三类递进任务：在HUST公开正常退化数据上获得可靠SOH估计，在自建不完整标定数据上维持低标定成本下的预测稳定性，并在自建异常电池数据上为故障识别提供轨迹级信息。基于此，本文构建由多尺度CNN、LSTM和物理一致性约束组成的SOH估计框架。


三个卷积分支的感受野具有不同的物理含义。kernel=3的短程分支更适合捕捉相邻循环之间的细微响应变化，例如充电时间或局部电压统计量的短期扰动；kernel=7的中程分支可覆盖若干连续循环，适合识别阶段性退化斜率变化；kernel=15的长程分支则更接近局部退化趋势估计，能够对缓慢变化的容量衰减模式形成更稳定响应。多尺度融合相当于让模型同时观察局部、阶段和趋势三个层面的证据，从而降低单一尺度对噪声或局部异常的敏感性。


在参数设计上，多尺度模块需要在表达能力和过拟合风险之间取得平衡。若每个分支通道数过大，模型可能在有限标签条件下记忆训练电池的局部模式；若通道数过小，则难以充分表达复杂退化特征。本文采用各分支相对轻量的通道配置，并在拼接后使用1×1卷积进行信息融合，使网络在保持计算复杂度可控的同时获得多粒度表示。


从工程部署角度看，多尺度卷积还具有计算结构清晰、并行度高和易于剪枝压缩等优点。相比Transformer或复杂注意力结构，多尺度CNN对序列长度的计算复杂度更可控，且卷积核大小与退化时间尺度之间具有较直观的解释关系。对于BMS边缘部署而言，模型不仅要追求离线测试误差，还需要考虑存储开销、推理时延和参数敏感性。因此，本文将多尺度卷积作为主线架构改进，而将注意力、Transformer和伪标签等更复杂策略放入扩展消融与边界讨论。



\subsection{4.3.3 LSTM时序建模与状态映射}

多尺度卷积模块输出的特征仍需在循环维度上进行长期状态建模。本文采用双层堆叠LSTM作为时序建模组件，其门控结构能够选择性保留历史状态信息，有助于捕获SOH随循环次数演化的长期依赖关系。LSTM最初用于缓解传统RNN中的梯度消失问题，适合处理具有长程依赖的序列建模任务[33]。在网络末端，采用全连接层与Sigmoid激活函数将高维隐状态映射为归一化SOH预测值，使输出自然位于[0,1]范围内。


LSTM在本文中的作用并不是替代物理模型，而是对多尺度卷积提取到的局部退化表征进行序列组织。卷积层更关注一定窗口内的局部模式，而LSTM通过门控机制维护跨循环的隐状态，使模型能够区分短期噪声与长期退化趋势。特别是在部分生命周期监督条件下，某些退化阶段缺少SOH标签，LSTM的状态传递有助于将有标签区间学到的退化模式迁移到相邻无标签区间。


Sigmoid输出层提供了基本的边界限制，但这并不等价于完整的物理一致性约束。Sigmoid只能保证SOH预测位于[0,1]范围内，不能保证预测序列随循环演化方向合理。因此，本文仍需在损失函数中引入软单调约束，用于控制预测轨迹的相对变化方向。


本文未将LSTM替换为更复杂的Transformer结构，主要基于两个考虑。第一，当前研究问题的核心不是长文本式全局依赖建模，而是中等长度生命周期序列中的退化趋势提取，LSTM已经能够提供足够的状态记忆能力。第二，Transformer虽然通过自注意力机制增强全局关联建模能力[34]，但在标签稀缺条件下更容易受到数据规模、正则化和超参数设置影响。因而本文将Transformer作为文献背景和可能扩展方向，而不是当前主线模型。



\subsection{4.3.4 基于软单调性的物理一致性弱监督}

锂离子电池在稳定退化阶段的SOH整体应呈非上升趋势。该现象来源于活性锂损失、活性材料损失、SEI膜增长和内阻增加等不可逆退化过程[23-24]。基于这一宏观先验，本文引入延迟激活的软单调约束。该约束并不要求所有相邻循环严格单调，而是允许一定容忍范围内的局部波动，仅对超过容忍阈值的非物理上升进行惩罚。


公式4-3  L\textbackslash{}\_mono = mean\textbackslash{}\_\textbackslash{}\{(i,t) in P\textbackslash{}\} max(0, yhat\textbackslash{}\_\textbackslash{}\{t+1\textbackslash{}\}\textbackslash{}textasciicircum\{\}\textbackslash{}\{(i)\textbackslash{}\} - yhat\textbackslash{}\_t\textbackslash{}textasciicircum\{\}\textbackslash{}\{(i)\textbackslash{}\} - epsilon)


其中，P表示满足约束激活条件的相邻循环对集合，epsilon为容忍因子。考虑到电池寿命早期可能存在容量活化、测试噪声或短暂平台波动，本文采用延迟激活策略，仅在循环索引超过给定阈值后施加单调约束。



\begin{figure}[H]
  \centering
  % \includegraphics[width=0.8\textwidth]{figures/placeholder}
  \caption{图4-5  锂离子电池生命周期早期容量回升现象示意图（图片占位）}
  \label{fig:}
\end{figure}

在完整标签条件下，软单调约束可视为对数据拟合项的辅助正则化；但在部分生命周期监督条件下，它具有更关键的机制意义。对于无标签循环，数据误差项不提供梯度，而软单调约束仍可基于相邻循环预测关系产生训练信号。因此，该约束不仅提高预测轨迹的物理合理性，也为无标签生命周期区间提供结构性弱监督。物理信息机器学习和PINN研究表明，将物理规律作为损失项或结构约束嵌入模型，有助于在小样本或不完全观测条件下提升泛化和可解释性[35-37]。


软单调约束采用延迟激活而非全生命周期强制约束，是因为实际电池在早期循环中可能出现短暂容量回升或平台波动。这类现象可能来源于电极润湿、活化过程、测试误差或温度微小扰动。若从第一个循环开始施加强单调约束，模型可能被迫忽略这些真实存在的早期变化，反而降低拟合质量。因此，本文将单调约束主要作用于稳定退化阶段，使物理先验与数据特性保持一致。


此外，本文采用“软”约束而非硬约束。硬约束要求每个相邻循环预测值严格不升高，虽然物理含义清晰，但在含噪测量数据和复杂工况下过于刚性。软约束通过容忍阈值允许小幅局部波动，只惩罚超过阈值的非物理上升。这种设计更符合工程数据的实际形态，也有助于避免物理约束与数据拟合目标发生剧烈冲突。


需要强调的是，本文的物理一致性约束属于轻量级物理先验，而非完整电化学模型求解。完整PINN通常需要明确的控制方程、边界条件或退化动力学方程；而实际BMS数据中，许多内部状态不可直接观测，完整电化学参数辨识成本较高。本文选择容量退化方向这一更稳定、可迁移的宏观先验，是为了在工程可实现性与物理解释性之间取得平衡。该策略不追求重建全部电化学机理，而是将最可靠的退化规律转化为训练约束。



\subsection{4.3.5 综合损失函数与评价指标}

模型训练目标由数据监督损失和物理一致性损失组成，如公式4-4所示。


公式4-4  L\textbackslash{}\_total = L\textbackslash{}\_data + lambda\textbackslash{}\_mono L\textbackslash{}\_mono


其中，L\textbackslash{}\_data仅在有标签循环上计算，如公式4-5所示。


在完整标定条件下，软单调约束可视为对数据拟合项的辅助正则化；但在不完整标定条件下，它具有更关键的机制意义。对于无标签循环，数据误差项不提供梯度，而软单调约束仍可基于相邻循环预测关系产生训练信号。因此，该约束不仅提高预测轨迹的物理合理性，也为无标签生命周期区间提供结构性弱监督。物理信息机器学习和PINN研究表明，将物理规律作为损失项或结构约束嵌入模型，有助于在小样本或不完全观测条件下提升泛化和可解释性[35-37]。


该损失设计体现了本文方法的核心逻辑：有标签循环通过数据损失学习准确的SOH数值映射，无标签循环通过物理一致性约束参与退化方向学习。评价指标包括RMSE、MAE、R²以及单调违反率。其中，RMSE和MAE衡量数值误差，R²衡量整体拟合能力，单调违反率衡量预测轨迹对退化单调性先验的遵循程度。


在优化过程中，数据误差项和物理一致性项承担不同职责。数据误差项决定模型输出的绝对SOH数值尺度，是保证预测精度的主要来源；物理一致性项决定预测序列的可行演化方向，是保证未标注区间稳定性的关键来源。若物理权重过小，无标签区间约束不足；若物理权重过大，模型可能牺牲局部拟合精度以满足过强的单调性要求。因此，物理权重需要通过验证集或消融实验确定，并在论文中明确讨论其适用边界。多任务学习中常通过不确定性或可学习权重平衡多个损失项[38]，但在本文低标签场景下，固定权重加消融分析更便于解释物理约束的真实贡献。


除常规误差指标外，本文引入单调违反率作为物理一致性评价指标。该指标统计同一电池预测序列中出现非物理上升的比例，可用于衡量模型输出是否符合容量退化的基本趋势。需要注意的是，单调违反率不能单独代表模型优劣；过低的违反率可能来自过度平滑，而非真实精度提升。因此，本文将MAE、RMSE、R²和单调违反率联合使用，以同时评价数值精度和物理可信度。


指标解释应避免“误差越低即完全更优”的单一判断。SOH估计面向的是安全运行和寿命管理，模型输出不仅要平均误差小，还要在退化方向、局部波动和未见电池上保持可信。尤其在部分生命周期监督场景中，某些模型可能在有标签区间误差较低，但在无标签区间产生明显非物理波动；另一些模型平均误差略高，却能给出更平滑和更符合退化规律的轨迹。因此，本文将数值误差、物理一致性和轨迹可解释性并列为评价对象。



\subsection{4.3.6 训练流程与实现要点}

MS-PI-CNNLSTM的训练流程围绕监督掩码展开。每个mini-batch中同时包含有标签循环和无标签循环。模型首先对所有循环特征进行前向传播，得到完整预测序列；随后，数据损失只在监督掩码为1的样本上计算，而物理一致性损失根据同一电池内部的循环顺序计算。这样既避免了无标签SOH对MSE造成污染，又保证无标签循环不会在训练中完全失效。


% 表4-3  MS-PI-CNNLSTM训练流程（伪代码形式）


\begin{table}[H]
  \centering
  \begin{tabular}{lll}
  \toprule
  步骤 & 操作 & 目的 \\
  \midrule
  1 & 读取电池级训练批次，包含特征、SOH标签、监督掩码和循环索引 & 保留完整生命周期特征与标签可见性信息 \\
  2 & 将特征输入多尺度CNN和LSTM，得到所有循环的SOH预测 & 对有标签和无标签循环统一前向建模 \\
  3 & 基于监督掩码计算有标签样本的MSE损失 & 学习特征到SOH的数值映射 \\
  4 & 按电池ID和循环索引构造相邻预测对，计算软单调损失 & 为无标签区间提供退化方向弱监督 \\
  5 & 合成总损失并反向传播，更新网络参数 & 联合优化预测精度与物理一致性 \\
  6 & 在验证集上监控MAE/RMSE和轨迹稳定性，执行早停 & 避免过拟合并选择稳定模型 \\
  \bottomrule
  \end{tabular}
  \caption{}
  \label{tab:}
\end{table}


实现上需要特别注意两个问题。第一，监督掩码必须只作用于数据误差项，不能影响模型对无标签循环的前向传播；否则模型无法利用完整生命周期特征轨迹。第二，物理一致性损失应基于同一电池内部的循环顺序计算，不能直接依赖mini-batch中的随机排列顺序。若训练时采用shuffle机制，需要在计算物理损失前根据电池ID和循环索引重新组织预测结果。


从可复现性角度看，本文建议在实验配置中明确报告随机种子、电池划分方式、监督比例、窗口长度、批大小、学习率调度策略、早停准则以及物理约束权重。对于部分生命周期监督任务，监督区间生成方式同样需要报告，因为不同标签覆盖位置可能对模型性能产生显著影响。公开电池数据研究已经越来越重视元数据、实验协议和评价口径的完整记录[16]，本文也应在最终论文中将这些信息写清楚。


训练流程还应避免两个容易造成结果偏差的细节。其一，标准化参数只能在训练集上拟合，再应用到验证集和测试集；若在全数据上拟合StandardScaler，会造成测试集统计信息泄露。其二，早停和超参数选择只能依赖验证集，测试集不应参与模型选择。对于SCI论文投稿而言，这些实验规范往往不会构成创新点，但会显著影响审稿人对结果可信度的判断。



\section{指标解释应避免“误差越低即完全更优”的单一判断。SOH估计面向的是安全运行和寿命管理，模型输出不仅要平均误差小，还要在退化方向、局部波动和未见电池上保持可信。在HUST正常退化数据上，重点考察数值精度和跨电池泛化；在自建不完整标定数据上，还需要关注无标定区间是否出现非物理波动；在异常电池数据上，则需要进一步考察预测轨迹偏差是否能够为故障识别提供有效信号。}

本节实验设计服务于“场景—机制—边界”的验证逻辑。首先验证部分生命周期监督是否构成独立挑战；其次分析多尺度架构是否提升有限标签下的特征表达；再次考察物理一致性约束何时能够为无标签区间提供有效弱监督；最后讨论注意力、不确定性估计、自适应权重和伪标签等扩展模块在该问题下的收益边界。由于当前部分实验仍需统一口径重跑，以下数值结果位置可在最终实验完成后回填。


% 表4-4  本章实验假设与验证路径


\begin{table}[H]
  \centering
  \begin{tabular}{llll}
  \toprule
  假设编号 & 科学问题 & 验证方式 & 支持该假设的预期现象 \\
  \midrule
  H1 & 部分生命周期监督会暴露完整监督下被掩盖的泛化问题 & 比较不同监督比例下的误差与轨迹稳定性 & 标签减少后部分模型误差增大或预测轨迹波动增强 \\
  H2 & 多尺度CNN能提升有限标签条件下的退化表征能力 & 比较CNN-LSTM与MS-CNN-LSTM & MS-CNN-LSTM在中低监督比例下表现更稳定 \\
  H3 & 软单调约束能为无标签区间提供结构性弱监督 & 比较无物理与软单调约束模型 & 低监督比例下物理约束降低非物理波动或改善轨迹稳定性 \\
  H4 & 复杂模块堆叠存在收益边界 & 比较M1/M2/M4/M5/M6等扩展模块 & 部分模块在低标签或校准不足时不稳定 \\
  \bottomrule
  \end{tabular}
  \caption{}
  \label{tab:}
\end{table}



\subsection{4.4.1 实验设置与对比方案}

所有实验均基于HUST数据集开展，严格遵循电池级别划分原则。模型输入为充电循环特征序列，输出为对应循环SOH估计值。部分监督实验设置多个监督比例r，用于模拟从完整标定到极度稀疏标定的不同工程场景。


从可复现性角度看，本文建议在实验配置中明确报告随机种子、电池划分方式、数据集来源、监督比例、窗口长度、批大小、学习率调度策略、早停准则以及物理约束权重。对于自建不完整标定任务，监督区间生成方式同样需要报告，因为不同标签覆盖位置可能对模型性能产生显著影响。公开电池数据研究已经越来越重视元数据、实验协议和评价口径的完整记录[16]，本文也应在最终论文中将这些信息写清楚。


% 表4-5  主线对比模型与配置说明


\begin{table}[H]
  \centering
  \begin{tabular}{llll}
  \toprule
  模型 & 架构 & 物理约束 & 作用 \\
  \midrule
  CNN-LSTM & 单尺度CNN+LSTM & 无 & 深度学习基线 \\
  PI-CNN-LSTM & 单尺度CNN+LSTM & 软单调 & 物理约束基线 \\
  MS-CNN-LSTM & 多尺度CNN+LSTM & 无 & 架构创新验证 \\
  MS-PI-CNNLSTM & 多尺度CNN+LSTM & 软单调 & 本文主线方法 \\
  \bottomrule
  \end{tabular}
  \caption{}
  \label{tab:}
\end{table}


本节实验设计服务于“公开基准—自建不完整标定—自建异常识别”的三阶段验证逻辑。首先，在HUST公开正常退化数据集上验证模型的基础SOH估计能力和跨电池泛化能力；其次，在自建不完整标定数据集上分析多尺度架构和物理一致性约束在低标定成本条件下的作用边界；最后，在自建异常电池数据集上考察SOH预测轨迹是否能够用于异常退化识别。由于当前部分实验仍需统一口径重跑，以下数值结果位置可在最终实验完成后回填。


实验结果报告时，应尽量采用均值与标准差同时呈现。部分生命周期监督下模型对随机种子和标签覆盖区间可能较敏感，仅报告单次运行结果容易夸大偶然收益。对于关键结论，建议至少比较多个随机种子下的平均表现，并结合代表性测试电池的轨迹图进行定性分析。不确定性估计相关研究表明，单模型点估计往往难以反映模型在分布外样本和低标签条件下的可信程度[39-40]，因此本文即使不把不确定性模块作为主线，也应在讨论中说明结果波动和置信解释的边界。


本章实验不宜只排列模型排行榜，而应围绕研究假设组织。完整监督实验用于建立基线，部分监督实验用于暴露标签稀缺带来的困难，架构消融用于解释多尺度表征的必要性，物理约束实验用于验证无标签区间弱监督的作用，鲁棒性实验用于说明工程扰动下的可用性。



\subsection{4.4.1 实验设置与三阶段验证方案}

本文实验分为三个层次。第一层为HUST公开数据集上的正常退化SOH估计实验，主要检验模型在公开标准数据上的预测精度和未见电池泛化能力。第二层为自建不完整标定数据集上的低标定成本验证，设置不同监督比例或不同标定覆盖区间，用于模拟工程维护中SOH标定不连续的情况。第三层为自建异常电池数据集上的故障识别扩展，利用SOH预测轨迹和异常分数判断电池是否偏离正常退化模式。


在撰写时，全监督基准部分应避免过度展开。该部分需要说明传统机器学习模型、纯时序模型和CNN-LSTM模型的大致性能层次，并确认本文模型在完整标签条件下不低于基本可用水平。真正体现论文价值的实验应放在后续部分监督场景，因为该场景更贴近实际BMS中SOH标定受限的应用条件。


全监督结果还可用于解释为什么仅在完整标签条件下评价模型是不充分的。若多个深度模型在r=1.0下差异较小，说明密集标签已经为模型提供了充分监督信号，使架构细节和物理先验的差异被弱化。此时，进一步比较部分监督条件下的表现更能反映模型对实际标定受限场景的适应能力。


% 表4-6  全监督条件下模型预测性能对比（待完整实验回填）


\begin{table}[H]
  \centering
  \begin{tabular}{lllll}
  \toprule
  模型 & RMSE(\%) & MAE(\%) & R² & 备注 \\
  \midrule
  XGBoost & 待回填 & 待回填 & 待回填 & 传统机器学习基线 \\
  LSTM & 待回填 & 待回填 & 待回填 & 时序模型基线 \\
  CNN-LSTM & 待回填 & 待回填 & 待回填 & 单尺度深度基线 \\
  MS-PI-CNNLSTM & 待回填 & 待回填 & 待回填 & 本文方法 \\
  \bottomrule
  \end{tabular}
  \caption{}
  \label{tab:}
\end{table}

除主线模型外，本文还将注意力机制、MC Dropout、速率连续性约束、自适应损失权重和伪标签训练作为扩展模块进行边界分析。这些模块不共同构成最终方法，而用于说明在正常退化估计、不完整标定和异常识别三个场景下，哪些增强策略稳定有效，哪些策略可能因标签稀缺、校准不足或异常模式不明确而失效。


实验结果报告时，应尽量采用均值与标准差同时呈现。不完整标定条件下模型对随机种子和标签覆盖区间可能较敏感，仅报告单次运行结果容易夸大偶然收益。对于关键结论，建议至少比较多个随机种子下的平均表现，并结合代表性测试电池的轨迹图进行定性分析。不确定性估计相关研究表明，单模型点估计往往难以反映模型在分布外样本和低标签条件下的可信程度[39-40]，因此本文即使不把不确定性模块作为主线，也应在讨论中说明结果波动和置信解释的边界。



\begin{figure}[H]
  \centering
  % \includegraphics[width=0.8\textwidth]{figures/placeholder}
  \caption{图4-7  全监督条件下模型性能指标对比（图片占位）}
  \label{fig:}
\end{figure}


\subsection{4.4.2 HUST公开正常退化数据集上的基准性能}

HUST公开数据集实验用于确认模型具备基本SOH估计能力，是本文三阶段验证路线的第一步。当SOH标签覆盖完整生命周期时，模型能够直接从密集标签中学习退化映射，多尺度结构和物理约束的边际收益可能被充分数据所掩盖。因此，本节结果主要作为后续自建不完整标定和异常识别实验的参照。


在撰写时，HUST基准部分应说明传统机器学习模型、纯时序模型和CNN-LSTM模型的大致性能层次，并确认本文模型在正常退化公开数据上达到可靠的基础预测水平。该部分不是为了直接证明低标定成本优势，而是为了建立可信的SOH估计基线：只有模型在正常退化数据上表现稳定，后续不完整标定和异常识别扩展才具有讨论基础。



\begin{table}[H]
  \centering
  \begin{tabular}{llll}
  \toprule
  架构 & 无物理约束 & 软单调约束 & 扩展物理/模块组合 \\
  \midrule
  CNN-LSTM & 待回填 & 待回填 & 待回填 \\
  MS-CNN-LSTM & 待回填 & 待回填 & 待回填 \\
  \bottomrule
  \end{tabular}
  \caption{}
  \label{tab:}
\end{table}

HUST结果还可用于解释为什么仅在完整标签条件下评价模型是不充分的。若多个深度模型在完整标定条件下差异较小，说明密集标签已经为模型提供了充分监督信号，使架构细节和物理先验的差异被弱化。此时，进一步比较自建不完整标定数据集上的表现更能反映模型对实际标定受限场景的适应能力。


若行方向，即CNN-LSTM到MS-CNN-LSTM的变化带来稳定改善，则说明多尺度表征是主要贡献；若列方向，即增加物理约束或扩展模块带来稳定改善，则说明约束强度是主要贡献。该矩阵能够避免将所有模块混合后难以解释贡献来源的问题。



\begin{figure}[H]
  \centering
  % \includegraphics[width=0.8\textwidth]{figures/placeholder}
  \caption{图4-8  架构×约束消融矩阵热力图（图片占位）}
  \label{fig:}
\end{figure}

从机制上看，多尺度架构解决的是输入表征不足问题，物理约束解决的是输出空间缺少结构限制问题。二者并非替代关系。若底层特征表征不足，物理约束只能限制输出形式，难以补偿输入信息缺失；若特征表征充分但标签稀缺，物理约束则可在无标签区间提供必要的退化方向引导。


建议在该小节中加入两类图形化结果。第一类是2×3矩阵热力图，用颜色直观表示不同架构和约束组合的MAE大小；第二类是相对改善柱状图，分别展示架构维度和约束维度的改善幅度。前者便于读者把握总体格局，后者便于突出“架构贡献大于约束堆叠”或“低标签下物理约束更有价值”等机制结论。


在解释架构消融时，应避免仅以“某个模型数值最低”作为结论。更稳健的写法是讨论趋势：多尺度结构是否在多个监督比例下保持更优或更稳定；软单调约束是否在标签稀缺时更有作用；Full Stack或复杂模块是否出现收益不稳定。这样的讨论更符合SCI论文对机制解释和可重复性的要求。


若最终实验显示多尺度结构在全监督条件下提升有限，但在中低监督比例下更稳定，这并不是负结果，而是非常符合本文问题设定的结果。因为多尺度结构的主要价值并非在密集标签下无限提高拟合能力，而是在标签减少时提高特征利用效率和退化趋势识别能力。该解释可与已有针对不完整工况和扩张卷积特征提取的研究形成呼应[28,32]，从而把实验现象转化为具有机制意义的论文结论。



\subsection{4.4.4 监督密度与物理约束作用边界分析}

为更贴近实际BMS应用场景，本文在不同监督比例下评估各模型性能。该实验用于回答：当SOH标签逐渐减少时，多尺度架构和物理一致性约束分别在何种条件下发挥作用。


% 表4-8  不同监督比例下各模型MAE对比（待完整实验回填）


\begin{table}[H]
  \centering
  \begin{tabular}{lllll}
  \toprule
  模型 & r=1.0 & r=0.7 & r=0.5 & r=0.3 \\
  \midrule
  CNN-LSTM & 待回填 & 待回填 & 待回填 & 待回填 \\
  PI-CNN-LSTM & 待回填 & 待回填 & 待回填 & 待回填 \\
  MS-CNN-LSTM & 待回填 & 待回填 & 待回填 & 待回填 \\
  MS-PI-CNNLSTM & 待回填 & 待回填 & 待回填 & 待回填 \\
  \bottomrule
  \end{tabular}
  \caption{}
  \label{tab:}
\end{table}



\begin{figure}[H]
  \centering
  % \includegraphics[width=0.8\textwidth]{figures/placeholder}
  \caption{图4-9  不同监督比例下各模型MAE变化趋势（图片占位）}
  \label{fig:}
\end{figure}

理论上，当监督比例较高时，数据损失已经提供充分退化趋势信息，物理约束可能只起到辅助正则作用，甚至在存在早期容量回升或局部波动时限制模型拟合能力。当监督比例较低时，无标签循环占比增加，物理一致性约束对这些样本提供退化方向弱监督，其作用应更加突出。该分析能够将物理约束的价值表述为具有场景边界的机制贡献，而不是无条件提升精度的经验技巧。


该部分还应讨论一个容易被忽略的现象：监督比例降低并不一定导致误差单调增大。若标签减少同时降低了过拟合风险，或物理约束在无标签区间发挥了更强正则化作用，低监督比例下的平均误差可能与高监督比例接近，甚至在某些随机划分中更低。因此，论文中不应简单将r作为“越大越好”的变量，而应结合监督覆盖位置、物理约束强度和模型容量共同解释。


从工程角度看，监督密度分析可为BMS标定策略提供参考。如果模型在中等监督比例下已经达到接近完整监督的性能，说明无需对每个循环进行高成本容量标定；如果极低监督比例下物理约束显著改善轨迹稳定性，说明轻量物理先验可降低标定成本。但若过强约束在高监督场景下损害精度，也提示工程部署时应根据标定条件动态选择约束强度。


4.4.4 自建不完整标定数据集上的监督密度与物理约束作用边界分析


为更贴近实际BMS应用场景，本文在自建不完整标定数据集上评估不同监督比例或不同标定覆盖方式下的模型性能。该实验用于回答：当SOH标签不再完整连续可得时，多尺度架构和物理一致性约束分别在何种条件下发挥作用。



\subsection{4.4.5 扩展模块的收益边界与复杂度分析}

为避免论文呈现为模块堆叠，本文将过去尝试的注意力机制、不确定性估计、速率连续性约束、自适应权重和伪标签训练统一纳入收益边界分析。该部分的价值在于说明部分生命周期监督下并非所有常见增强策略都稳定有效。


% 表4-9  扩展模块的机制定位与论文处理方式


\begin{table}[H]
  \centering
  \begin{tabular}{llll}
  \toprule
  模块 & 设计初衷 & 理论风险 & 论文定位 \\
  \midrule
  M1 注意力 & 突出窗口内关键循环 & 低标签下可能过拟合局部有标签片段 & 消融项 \\
  M2 MC Dropout & 估计预测不确定性 & 方差不一定等价于真实误差 & 辅助分析 \\
  M4 速率连续性 & 抑制预测斜率突变 & 可能压制真实阶段性加速退化 & 物理扩展约束 \\
  M5 自适应权重 & 自动平衡损失项 & 低标签下权重可能漂移 & 负例或限制讨论 \\
  M6 伪标签 & 补充无标签样本监督 & 校准不足时高置信伪标签可能错误 & 限制讨论 \\
  \bottomrule
  \end{tabular}
  \caption{}
  \label{tab:}
\end{table}


上述模块均具有合理动机，但其有效性依赖更强训练信号或可靠不确定性校准。在SOH标签稀缺时，复杂模块可能放大局部偏差或引入额外不稳定性。相比之下，多尺度CNN和软单调约束分别对应更强表征和更明确先验，与部分监督问题的核心矛盾直接匹配，因此更适合作为最终论文主线。


注意力机制的风险在于，窗口内注意力权重可能过度集中于少数有标签或局部误差较小的循环，从而忽略整体退化趋势。虽然注意力和Transformer结构已经在序列建模中展示出强大的全局依赖建模能力[34]，电池SOH研究中也有混合注意力网络取得较好效果[41]，但在低标签条件下，注意力权重本身未必等价于可靠的物理解释。MC Dropout虽然可以提供不确定性估计[39]，但其预测方差不一定与真实误差严格相关，尤其在低标签条件下可能出现校准不足。伪标签方法则进一步依赖不确定性质量，一旦高置信样本并不高质量，错误伪标签会被模型反复学习并放大。


该分析还可以自然引出模型选择策略：当BMS系统能够获得较密集的SOH标定时，可优先采用纯多尺度架构以最大化数据拟合能力；当标定极少或生命周期覆盖不完整时，应启用软单调物理约束以降低无标签区间的预测自由度；当工况变化剧烈、特征异常明显或电池存在潜在故障时，则需要结合异常识别和轨迹偏差分析进一步约束模型输出。


这种边界分析本身具有论文价值。它说明在标定受限条件下，模型设计不能只追求结构复杂度，而应优先考虑模块与问题机制是否匹配。多尺度CNN直接对应退化特征的多粒度表达，软单调约束直接对应无标签区间的退化方向约束，因此比泛化的复杂模块更容易形成稳定、可解释和可部署的方案。


在硕士论文写作中，扩展模块不一定都要被写成“正向创新”。如果实验发现某些模块收益不稳定，可以将其放入“边界与限制”讨论：例如MC Dropout能够提供预测区间，但校准质量需要额外验证[39-40]；半监督伪标签能够利用无标签数据，但会引入伪标签错误累积风险[30]；注意力机制能够突出关键片段，但注意力权重可能受标签覆盖位置影响。这种写法比简单删除负结果更有说服力，也能展示研究过程的系统性。



\subsection{为避免论文呈现为模块堆叠，本文将过去尝试的注意力机制、不确定性估计、速率连续性约束、自适应权重和伪标签训练统一纳入收益边界分析。该部分的价值在于说明三阶段验证中并非所有常见增强策略都稳定有效：在HUST正常退化数据上有效的模块，未必能在不完整标定或异常识别场景中继续产生稳定收益。}

为进一步验证模型在未见电池上的泛化能力，本文选取测试集中若干代表性电池绘制SOH预测轨迹和误差演化曲线。重点观察模型是否能够跟随整体退化趋势，是否在无标签区间出现非物理回弹或局部突降，以及误差是否存在系统性偏移。



\begin{figure}[H]
  \centering
  % \includegraphics[width=0.8\textwidth]{figures/placeholder}
  \caption{图4-10  未见电池SOH预测轨迹对比（图片占位）}
  \label{fig:}
\end{figure}


\begin{figure}[H]
  \centering
  % \includegraphics[width=0.8\textwidth]{figures/placeholder}
  \caption{图4-11  未见电池预测误差演化曲线（图片占位）}
  \label{fig:}
\end{figure}

除干净测试集外，工程部署还需要考虑特征缺失、传感器噪声和漂移等扰动。本文建议将鲁棒性验证作为工程意义补充，比较基线模型与MS-PI-CNNLSTM在不同扰动强度下的绝对误差和误差增量。已有针对不完整工况、无标签充电数据和海事电池实测数据的研究均说明，真实应用中的SOH估计往往面对不完整观测、弱标签和运行条件变化[28,30-31]。若多尺度架构在特征缺失下保持更低误差，可解释为多粒度特征带来了更强冗余性；若噪声或漂移下灵敏度未显著增加，则说明模型复杂度没有引入额外脆弱性。


未见电池轨迹图在论文中非常重要。数值指标能够概括总体误差，但无法展示模型是否在局部区间出现不合理波动。建议选取寿命长度、退化斜率和局部噪声特征不同的若干测试电池，分别绘制真实SOH、CNN-LSTM预测、MS-CNN-LSTM预测和MS-PI-CNNLSTM预测曲线。这样可以直观说明多尺度结构改善趋势跟踪能力，物理约束抑制非物理回弹或过度波动。


误差演化曲线则用于分析模型偏差来源。如果某模型在退化早期误差较小但后期误差逐渐扩大，说明其对加速退化阶段泛化不足；如果误差呈现高频振荡，说明模型对局部噪声过敏；如果误差整体偏正或偏负，则说明模型存在系统性高估或低估。通过轨迹图和误差图结合，可以比单一MAE更充分地解释模型行为。


% 表4-10  工程鲁棒性验证设计


\begin{table}[H]
  \centering
  \begin{tabular}{llll}
  \toprule
  扰动类型 & 扰动方式 & 评价指标 & 工程含义 \\
  \midrule
  特征缺失 & 随机屏蔽若干输入特征 & MAE与ΔMAE & 模拟传感器短时失效或数据缺测 \\
  高斯噪声 & 输入特征加入不同强度噪声 & MAE与ΔMAE & 模拟测量噪声 \\
  传感器漂移 & 输入特征加入缓慢偏移 & MAE与ΔMAE & 模拟传感器长期偏置 \\
  \bottomrule
  \end{tabular}
  \caption{}
  \label{tab:}
\end{table}

4.4.6 未见电池泛化与跨场景工程鲁棒性验证


鲁棒性实验的写法应同时报告绝对误差和误差增量。绝对误差反映模型在扰动条件下的实际可用性，误差增量反映模型对扰动的敏感程度。某些情况下，本文模型的误差增量可能与基线相近，但由于干净条件下误差更低，扰动后的绝对误差仍然更优。这种结果在工程上同样有意义，因为BMS最终关心的是扰动环境下的估计误差是否仍处于可接受范围内。


如果后续加入异常检测扩展实验，建议将其作为工程应用拓展而非本章核心贡献。异常检测可利用预测轨迹偏离、局部下降速率异常或特征分布偏离等信号，但其有效性依赖故障注入方式和评分函数设计。为了不稀释SOH估计主线，异常检测更适合在第四章末尾作为附加验证，或在后续章节/论文扩展版本中单独展开。


工程鲁棒性部分的写作重点应放在“扰动后模型行为是否仍可解释”。例如，在特征缺失实验中，如果长程分支仍能维持退化趋势，则说明多尺度表示具有一定信息冗余；在噪声实验中，如果物理约束降低了预测高频振荡，则说明其对输出轨迹具有稳定作用；在漂移实验中，如果模型出现系统性偏差，则应讨论标准化策略、在线校准或域适应方法的必要性。这样的分析能够使鲁棒性实验从简单附加测试上升为工程可部署性论证。



\subsection{4.4.7 基于SOH预测轨迹的零训练代价异常检测扩展}

在完成SOH估计主线实验后，本文进一步考察已训练模型在异常退化感知中的扩展价值。该实验的核心问题是：在不引入真实失效样本、不重新训练异常检测模型、不改变SOH估计网络权重的条件下，能否仅依赖模型推理阶段产生的SOH预测轨迹和输入特征统计量，对异常退化进行早期提示。本文将该思路称为零训练代价异常检测，其意义不在于替代专用故障诊断模型，而在于验证SOH估计模型是否具备一定的工程复用能力。


该扩展实验与本章主线之间的关系需要明确界定。MS-PI-CNNLSTM的核心贡献仍然是面向部分生命周期监督的SOH估计，即在SOH标签不完整条件下利用多尺度特征表征和物理一致性弱监督提高未见电池预测稳定性。异常检测实验只是基于同一预测模型的后验应用：模型权重固定，输入正常或异常注入后的特征序列，输出SOH预测轨迹，再根据轨迹偏离程度构造异常分数。因此，该实验应作为工程应用拓展而非方法主贡献呈现。


异常场景构造采用self-trajectory策略，即在同一电池生命周期内部将未来老化状态提前注入到故障点之后，而不是从其他电池混入末期特征。前期跨电池注入实验表明，不同电池之间的个体差异可能掩盖老化时间差异，导致模型对异常注入不敏感；相反，同电池轨迹重排既保留个体特征，又能模拟提前老化、加速退化或突发容量损失，更适合评价SOH预测轨迹对异常退化的响应能力。本文统一将故障注入位置设为50\textbackslash{}\%生命周期处，并设置mild、moderate和severe三级严重度。


% 表4-11  零训练代价异常检测中的五类退化场景设计


\begin{table}[H]
  \centering
  \begin{tabular}{llll}
  \toprule
  场景 & 物理含义 & 注入方式 & 预测轨迹表现 \\
  \midrule
  A1 突发老化 & 冲击、过充/过放或短路导致不可逆容量骤降 & 跳过中间健康片段，拼接同电池更低SOH状态 & SOH预测出现明显台阶下跌 \\
  A2 容量拐点 & 活性锂损失累积后退化速率突然增大 & 故障点后按更大步长跳采样未来状态 & 预测曲线斜率在拐点后变陡 \\
  A3 簇内不均衡 & 电池组短板单体加速老化导致整体SOH持续偏低 & 线性混合同电池未来特征，混合比例逐步增大 & 预测轨迹相对正常轨迹持续下偏 \\
  A4 析锂 & 低温或大倍率充电导致活性锂损失并引发加速衰退 & 小台阶下跌与跳采样加速组合 & 台阶下跌后继续加速退化 \\
  A5 内阻增长 & SEI膜增厚和接触电阻增大导致缓慢加速老化 & 按二次增长比例混入未来状态 & 早期偏移弱，后期斜率逐渐增大 \\
  \bottomrule
  \end{tabular}
  \caption{}
  \label{tab:}
\end{table}


异常检测扩展实验应作为第三阶段应用验证而非本章核心SOH估计贡献。异常检测可利用预测轨迹偏离、局部下降速率异常或特征分布偏离等信号，但其有效性依赖故障样本来源、异常注入方式和评分函数设计。为了不稀释SOH估计主线，异常识别更适合在HUST基准和自建不完整标定实验之后展开，作为模型工程复用能力的补充说明。


上述五类场景并非真实故障标签数据，而是基于正常退化轨迹构造的可控合成异常。其优点是可以在缺少真实失效电池的条件下评估检测信号是否具备基本区分能力；其局限是不能完全覆盖实际电池故障中的热失控、连接故障、传感器异常或复杂工况耦合。因此，本文仅将其作为工程可扩展性验证，结论表述应保持谨慎。


4.4.7 基于自建异常电池数据集的SOH轨迹异常识别扩展


在完成HUST正常退化数据和自建不完整标定数据上的SOH估计实验后，本文进一步考察已训练模型在异常退化感知中的扩展价值。该实验的核心问题是：在不重新训练复杂故障诊断模型的条件下，能否利用SOH预测模型在异常电池数据上的推理轨迹，对电池异常退化或潜在故障进行早期提示。本文将该思路定位为SOH估计模型的工程复用验证，其意义不在于替代专用故障诊断模型，而在于说明健康状态估计结果能否进一步服务于安全预警。


该扩展实验与本章主线之间的关系需要明确界定。MS-PI-CNNLSTM的核心贡献首先是在HUST公开正常退化数据上获得可靠SOH估计，其次是在自建不完整标定数据上验证低标定成本条件下的适用性；异常识别实验则是基于同一预测模型的后验应用。模型权重固定，输入正常或异常电池特征序列，输出SOH预测轨迹，再根据轨迹偏离程度构造异常分数。因此，该实验应作为工程应用拓展而非方法主贡献呈现。


异常场景构造可采用自建异常电池实测数据，也可在真实异常样本不足时采用self-trajectory策略进行可控注入。self-trajectory方法在同一电池生命周期内部将未来老化状态提前注入到故障点之后，而不是从其他电池混入末期特征。前期跨电池注入实验表明，不同电池之间的个体差异可能掩盖老化时间差异，导致模型对异常注入不敏感；相反，同电池轨迹重排既保留个体特征，又能模拟提前老化、加速退化或突发容量损失，更适合评价SOH预测轨迹对异常退化的响应能力。



\begin{table}[H]
  \centering
  \begin{tabular}{llll}
  \toprule
  信号 & 计算依据 & 主要作用 & 局限性 \\
  \midrule
  input\_zscore & 输入特征相对训练集统计量的最大z-score & 识别特征分布偏离 & 正常寿命末期也可能偏高，易造成假阳性 \\
  mono\_violation & 相邻SOH预测的正向上升量 & 识别违反单调退化的异常波动 & 对加速退化和渐进偏移不敏感 \\
  rate\_anomaly & 当前退化速率相对局部均值的z-score & 识别退化速度突变 & 窗口统计噪声较大 \\
  drop\_anomaly & 当前下跌幅度相对局部基准的超限程度 & 识别异常骤降 & 对缓慢漂移类异常响应弱 \\
  trajectory\_deviation & 当前预测值与局部趋势外推值的差异 & 识别突变和累积轨迹偏离 & 窗口长度影响灵敏度与检测延迟 \\
  \bottomrule
  \end{tabular}
  \caption{}
  \label{tab:}
\end{table}


从论文写作逻辑看，本节应重点验证trajectory\textbackslash{}\_deviation是否能够成为最稳定的异常感知信号。input\textbackslash{}\_zscore虽然能够反映输入特征偏离训练分布的程度，但正常电池寿命后期也可能自然偏高，存在系统性假阳性风险；mono\textbackslash{}\_violation、rate\textbackslash{}\_anomaly和drop\textbackslash{}\_anomaly分别关注预测上升、速率突变和异常下跌，但对缓慢累积型异常可能不够敏感。原始combined信号若采用多个归一化分数的最大值，则容易被单一高噪声信号主导。因此，最终实验回填时建议优先报告trajectory\textbackslash{}\_deviation及其与其他信号的对比，而不是简单叠加所有信号。


% 表4-13  Exp-12中trajectory\_deviation信号在五类异常场景下的检测结果（待最终实验回填）


\begin{table}[H]
  \centering
  \begin{tabular}{llll}
  \toprule
  异常场景 & SOH变化形态 & AUC & 结果解释 \\
  \midrule
  A1 突发老化 & 明显台阶下跌 & 待回填 & 突变会造成局部趋势外推误差迅速增大，理论上应较易检测 \\
  A2 容量拐点 & 斜率加速下滑 & 待回填 & 拐点后趋势变化应能引起轨迹偏差增加，但通常弱于突发型异常 \\
  A3 簇内不均衡 & 持续线性偏低 & 待回填 & 偏移逐步累积，检测效果取决于异常幅度与模型随机误差的相对大小 \\
  A4 析锂 & 小台阶下跌并加速衰退 & 待回填 & 台阶和加速共同放大轨迹偏差，理论上应具有较清晰响应 \\
  A5 内阻增长 & 斜率缓慢二次增大 & 待回填 & 早期异常信号较弱，通常需要更长窗口或累计偏差后才能识别 \\
  \bottomrule
  \end{tabular}
  \caption{}
  \label{tab:}
\end{table}

上述五类场景可由自建异常电池实验获得，也可在真实失效样本不足时作为可控合成异常进行补充。其优点是可以在缺少大量真实失效电池的条件下评估检测信号是否具备基本区分能力；其局限是不能完全覆盖实际电池故障中的热失控、连接故障、传感器异常或复杂工况耦合。因此，本文仅将其作为工程可扩展性验证，结论表述应保持谨慎。



\begin{figure}[H]
  \centering
  % \includegraphics[width=0.8\textwidth]{figures/placeholder}
  \caption{图4-13  代表性测试电池的trajectory\_deviation异常分数时序曲线及报警阈值（图片占位）}
  \label{fig:}
\end{figure}

报警阈值采用正常时段异常分数的第95百分位确定，对应误报率约束FPR=5\textbackslash{}\%。若进一步采用连续N次超限确认机制，则可以显著降低偶然噪声导致的误报警。例如，单点误报率为5\textbackslash{}\%时，连续3次误报的理论概率约为0.05\textbackslash{}textasciicircum\{\}3，即0.0125\textbackslash{}\%，代价是报警时间可能延迟若干循环。该机制符合BMS实际部署中宁可牺牲少量响应速度以换取更低误报率的工程需求。


公式4-7  theta = percentile(\textbackslash{}\{s\textbackslash{}\_traj(t), t < t\textbackslash{}\_fault\textbackslash{}\}, 95),  alarm\textbackslash{}\_t = I[s\textbackslash{}\_traj(t)>theta]


待最终实验数据补充后，本节可围绕三个结论展开。第一，不额外训练异常检测模型，仅复用SOH预测轨迹，有望对突发老化、析锂等具有明显轨迹扰动的异常形成早期提示。第二，异常检测能力预期主要来自局部趋势外推偏差，而不是某个特定网络结构本身；若不同SOH估计模型之间AUC差异较小，则说明检测信号设计比模型堆叠更关键。第三，渐进型异常预计仍然较难检测，因为其异常幅度可能与模型预测随机误差处于相近量级，零训练代价方法更适合作为风险提示工具，而不能替代专门故障诊断模型。


因此，本文将Exp-12定位为工程扩展验证而非核心贡献。它说明SOH估计模型不仅可以输出健康状态数值，还能通过预测轨迹的局部一致性为异常退化提供早期提示；同时也暴露出合成异常、窗口长度选择、渐进型故障灵敏度和真实故障验证不足等限制。后续工作可围绕自适应窗口、分电池个性化基线、异常分数校准以及真实失效数据验证进一步展开。



\subsection{4.4.8 本章实验讨论}

综合上述实验设计，本文希望形成的主要结论不是某个单一模型在所有指标上绝对最优，而是明确不同机制在部分生命周期监督任务中的作用边界。多尺度CNN主要解决有限标签下的特征表达问题；LSTM负责跨循环状态演化；软单调约束为无标签区间提供物理方向引导；扩展模块需要在不确定性校准、标签覆盖和训练稳定性足够可靠时才可能进一步发挥作用。


这种讨论方式有助于增强论文说服力。SCI论文通常不仅要求报告性能提升，还要求说明为什么提升、何时提升以及何时不提升。对于控制工程和电池管理系统领域而言，模型的可解释性、部署代价和失效边界往往与平均误差同样重要。因此，将负结果和不稳定结果整理为边界分析，反而能够使论文显得更加严谨。


最终定稿时，建议将所有实验结果统一到同一套数据划分、同一组随机种子和同一评价脚本下重新生成。若某些历史结果来自不同训练协议或不同模型配置，应避免直接放入同一张主表。可以将其作为方法探索过程保留在附录或讨论中，而主文只呈现口径一致、可重复的关键结果。


从全文逻辑看，本章最终应形成四层递进关系：第一层说明SOH标签稀缺是真实工程痛点；第二层提出Label Masking将该痛点转化为可实验验证的部分生命周期监督问题；第三层用多尺度CNN-LSTM和软单调约束分别解决表征不足与无标签区间缺乏约束的问题；第四层通过监督密度、消融矩阵和鲁棒性实验说明该方法何时有效、为什么有效以及在哪些条件下仍有局限。该逻辑比单纯罗列模块更符合硕士论文和SCI论文的论证要求。



\section{4.5 本章小结}

本章围绕实际BMS中SOH标定不完整、生命周期标签覆盖有限的工程问题，提出一种面向部分生命周期监督的多尺度物理一致性SOH估计方法MS-PI-CNNLSTM。与传统完整生命周期监督假设不同，本文将运行特征连续可得而SOH标签稀疏可得的场景显式建模为Label Masking问题。


在方法设计方面，本文首先通过多尺度并行卷积提取短程、中程和长程退化特征，增强有限标签条件下的表征能力；随后利用LSTM建模跨循环SOH状态演化；最后通过延迟激活的软单调约束为无标签生命周期区间提供结构性弱监督。该设计使有标签循环通过数据损失学习数值映射，无标签循环通过物理一致性约束参与退化方向学习。


在实验组织方面，本章建议采用“场景—机制—边界”的验证逻辑：通过不同监督比例评估部分生命周期监督挑战，通过架构×约束矩阵解耦多尺度表征和物理约束贡献，通过扩展模块消融说明复杂模块并非越多越好，通过未见电池和扰动实验补充工程鲁棒性分析。


综合上述实验设计，本文希望形成的主要结论不是某个单一模型在所有指标上绝对最优，而是明确同一模型在三个场景中的作用边界。HUST公开数据集用于证明模型在正常退化条件下具备可靠SOH估计能力；自建不完整标定数据集用于检验多尺度表征和物理一致性约束在低标定成本下的价值；自建异常电池数据集用于考察SOH预测轨迹能否为故障识别提供早期提示。这样的实验组织比单一数据集上的性能比较更能体现工程意义。


本章仍有若干需要在最终实验后进一步完善的内容。首先，表格中的性能指标需要基于统一协议回填，并报告均值和标准差；其次，图4-1至图4-11需要根据最终数据重新绘制，确保图表与正文结论一致；再次，参考文献需要在投稿前按目标期刊格式统一大小写、缩写和DOI。完成上述工作后，本章可形成一条较完整的SCI论文叙事链：工程问题来自SOH标签不完整，方法设计对应多尺度表征和物理弱监督，实验分析揭示机制贡献和复杂度边界，鲁棒性验证体现工程应用意义。




\section{从全文逻辑看，本章最终应形成四层递进关系：第一层说明SOH估计是BMS寿命管理的基础任务；第二层在HUST公开正常退化数据集上建立可信的基础估计结果；第三层在自建不完整标定数据集上验证低标定成本场景中的预测稳定性；第四层在自建异常电池数据集上探索SOH预测轨迹对故障识别的辅助价值。该逻辑比一开始直接进入不完整监督更容易被读者接受，也更符合从公开基准到工程扩展的论文叙事方式。}

\begin{itemize}
  \item Severson, K. A., Attia, P. M., Jin, N., Perkins, N., Jiang, B., Yang, Z., et al. Data-driven prediction of battery cycle life before capacity degradation. Nature Energy, 2019, 4: 383-391.
\end{itemize}
\begin{itemize}
  \item 本章围绕锂离子电池SOH估计与工程应用扩展问题，提出一种面向多场景验证的多尺度物理一致性SOH估计方法MS-PI-CNNLSTM。与直接从不完整标定场景切入不同，本文首先在HUST公开正常退化数据集上验证模型的基础估计能力，再在自建不完整标定数据集上考察低标定成本条件下的适用性，最后在自建异常电池数据集上探索SOH预测轨迹的故障识别价值。
\end{itemize}
\begin{itemize}
  \item Ng, M. F., Zhao, J., Yan, Q., Conduit, G. J., \& Seh, Z. W. Predicting the state of charge and health of batteries using data-driven machine learning. Nature Machine Intelligence, 2020, 2: 161-170.
\end{itemize}
\begin{itemize}
  \item 在实验组织方面，本章建议采用“公开基准—不完整标定—异常识别”的验证逻辑：通过HUST公开数据集评估正常退化SOH估计能力，通过自建不完整标定数据集分析监督密度和物理约束作用边界，通过自建异常电池数据集补充工程安全预警能力分析。该组织方式能够使方法贡献、工程痛点和应用扩展之间形成更清晰的递进关系。
\end{itemize}
\begin{itemize}
  \item 总体而言，MS-PI-CNNLSTM并非依赖复杂模块堆叠，而是围绕SOH估计从公开基准到工程应用的递进需求构建可解释建模路线：以多尺度卷积提升正常退化特征表征能力，以LSTM连接跨循环状态演化，以软单调约束提高预测轨迹物理可信度，并在不完整标定与异常识别场景中进一步验证其工程价值。该路线兼具问题针对性、机制解释性与可实现性，可为锂离子电池SOH估计和故障风险提示提供参考。
\end{itemize}
\begin{itemize}
  \item Yang, N., Song, Z., Hofmann, H., \& Sun, J. Robust State of Health estimation of lithium-ion batteries using convolutional neural network and random forest. Journal of Energy Storage, 2022, 48: 103857.
\end{itemize}
\begin{itemize}
  \item Liu, K., Kang, L., \& Xie, D. Online State of Health Estimation of Lithium-Ion Batteries Based on Charging Process and Long Short-Term Memory Recurrent Neural Network. Batteries, 2023, 9(2): 94.
\end{itemize}
\begin{itemize}
  \item Zhu, Z., Yang, Q., Liu, X., \& Gao, D. Attention-based CNN-BiLSTM for SOH and RUL estimation of lithium-ion batteries. Journal of Algorithms \& Computational Technology, 2022, 16: 1-11.
\end{itemize}
\begin{itemize}
  \item Fan, Y., Li, Y., Zhao, J., Wang, L., Yan, C., Wu, X., et al. Online State-of-Health Estimation for Fast-Charging Lithium-Ion Batteries Based on a Transformer-Long Short-Term Memory Neural Network. Batteries, 2023, 9(11): 539.
\end{itemize}
\begin{itemize}
  \item  Lu, J., Xiong, R., Tian, J., Wang, C., \& Sun, F. Deep learning to estimate lithium-ion battery state of health without additional degradation experiments. Nature Communications, 2023, 14: 2760.
\end{itemize}
\begin{itemize}
  \item  Ismail Fawaz, H., Lucas, B., Forestier, G., Pelletier, C., Schmidt, D. F., Weber, J., Webb, G. I., et al. InceptionTime: Finding AlexNet for time series classification. Data Mining and Knowledge Discovery, 2020, 34(6): 1936-1962.
\end{itemize}
\begin{itemize}
  \item  毛百海, 覃吴, 肖显斌, 郑宗明. 基于LSTM\&GRU-Attention多联合模型的锂离子电池SOH估计. 储能科学与技术, 2023, 12(11): 3519-3527. DOI: 10.19799/j.cnki.2095-4239.2023.0514.
\end{itemize}
\begin{itemize}
  \item  戴彦文, 于艾清. 基于健康特征参数的CNN-LSTM\&GRU组合锂电池SOH估计. 储能科学与技术, 2022, 11(5): 1641-1649. DOI: 10.19799/j.cnki.2095-4239.2021.0623.
\end{itemize}
\begin{itemize}
  \item  王英楷, 张红, 王星辉. 基于1DCNN-LSTM的锂离子电池SOH预测. 储能科学与技术, 2022, 11(1): 240-245. DOI: 10.19799/j.cnki.2095-4239.2021.0250.
\end{itemize}
\begin{itemize}
  \item  张吉昂, 王萍, 程泽. 基于充电电压片段和融合方法的锂离子电池SOC-SOH-RUL联合估计. 电网技术, 2021. DOI: 10.13335/j.1000-3673.pst.2021.0888.
\end{itemize}
\begin{itemize}
  \item  Ward, L., Babinec, S., Dufek, E. J., Howey, D. A., Viswanathan, V., Aykol, M., et al. Principles of the Battery Data Genome. Joule, 2022, 6(10): 2253-2271.
\end{itemize}
\begin{itemize}
  \item  Yuan, Y., Ma, G., \& Xu, S. The Dataset for: Real-time personalized health status prediction of lithium-ion batteries using deep transfer learning. Mendeley Data, 2022, V2, DOI: 10.17632/nsc7hnsg4s.2.
\end{itemize}
\begin{itemize}
  \item  Wang, S., Sankarasubramanian, S., Chehade, A., \& Zhang, Y. Physics-informed neural network for lithium-ion battery degradation stable modeling and prognosis. Nature Communications, 2024, 15: 4332.
\end{itemize}
\begin{itemize}
  \item  Ye, J., Xie, Q., Lin, M., \& Wu, J. A method for estimating the state of health of lithium-ion batteries based on physics-informed neural network. Energy, 2024, 294: 130828.
\end{itemize}
\begin{itemize}
  \item  Lin, C., Tuo, X., Wu, L., Zhang, G., Lyu, Z., \& Zeng, X. Physics-informed machine learning for accurate SOH estimation of lithium-ion batteries considering various temperatures and operating conditions. Energy, 2025, 318: 134937.
\end{itemize}
\begin{itemize}
  \item  Lin, C., Wu, L., Tuo, X., Liu, C., Zhang, W., Huang, Z., \& Zhang, G. A lightweight two-stage physics-informed neural network for SOH estimation of lithium-ion batteries with different chemistries. Journal of Energy Chemistry, 2025, 105: 261-279.
\end{itemize}
\begin{itemize}
  \item  Ma, G., Xu, S., Jiang, B., Cheng, C., Yang, X., Shen, Y., et al. Real-time personalized health status prediction of lithium-ion batteries using deep transfer learning. Energy \& Environmental Science, 2022, 15: 4083-4094.
\end{itemize}
\begin{itemize}
  \item  Che, Y., Hu, X., Lin, X., Guo, J., \& Teodorescu, R. Health prognostics for lithium-ion batteries: mechanisms, methods, and prospects. Energy \& Environmental Science, 2023, 16: 338-371.
\end{itemize}
\begin{itemize}
  \item  Birkl, C. R., Roberts, M. R., McTurk, E., Bruce, P. G., \& Howey, D. A. Degradation diagnostics for lithium ion cells. Journal of Power Sources, 2017, 341: 373-386.
\end{itemize}
\begin{itemize}
  \item  Dubarry, M., \& Anseán, D. Best practices for incremental capacity analysis. Frontiers in Energy Research, 2022, 10: 1023555.
\end{itemize}
\begin{itemize}
  \item  Lyu, Z., Gao, R., \& Li, X. A partial charging curve-based data-fusion-model method for capacity estimation of Li-Ion battery. Journal of Power Sources, 2021, 483: 229131.
\end{itemize}
\begin{itemize}
  \item  Chen, Z., Sun, M., Shu, X., Xiao, R., Shen, J., et al. State of health estimation for lithium-ion batteries based on temperature prediction and gated recurrent unit neural network. Journal of Power Sources, 2022, 521: 230892.
\end{itemize}
\begin{itemize}
  \item  Zhang, Y., Wang, Y., Xia, Y., \& Chen, W. A deep learning approach to estimate the state of health of lithium-ion batteries under varied and incomplete working conditions. Journal of Energy Storage, 2023, 58: 106323.
\end{itemize}
\begin{itemize}
  \item  Li, Y., Sheng, H., Cheng, Y., Stroe, D.-I., \& Teodorescu, R. State-of-health estimation of lithium-ion batteries based on semi-supervised transfer component analysis. Applied Energy, 2020, 277: 115504.
\end{itemize}
\begin{itemize}
  \item  Lin, C., Xu, J., \& Mei, X. Improving state-of-health estimation for lithium-ion batteries via unlabeled charging data. Energy Storage Materials, 2023, 54: 85-97.
\end{itemize}
\begin{itemize}
  \item  De Bin, R., Furlan, L., Gianoli, P., \& Panza, A. A novel semi-supervised learning approach for State of Health monitoring of maritime lithium-ion batteries. Journal of Power Sources, 2023, 556: 232429.
\end{itemize}
\begin{itemize}
  \item  Bai, S., Kolter, J. Z., \& Koltun, V. An empirical evaluation of generic convolutional and recurrent networks for sequence modeling. arXiv preprint arXiv:1803.01271, 2018.
\end{itemize}
\begin{itemize}
  \item  Hochreiter, S., \& Schmidhuber, J. Long short-term memory. Neural Computation, 1997, 9(8): 1735-1780.
\end{itemize}
\begin{itemize}
  \item  Vaswani, A., Shazeer, N., Parmar, N., Uszkoreit, J., Jones, L., Gomez, A. N., et al. Attention is all you need. Advances in Neural Information Processing Systems, 2017, 30.
\end{itemize}
\begin{itemize}
  \item  Raissi, M., Perdikaris, P., \& Karniadakis, G. E. Physics-informed neural networks: A deep learning framework for solving forward and inverse problems involving nonlinear partial differential equations. Journal of Computational Physics, 2019, 378: 686-707.
\end{itemize}
\begin{itemize}
  \item  Karniadakis, G. E., Kevrekidis, I. G., Lu, L., Perdikaris, P., Wang, S., \& Yang, L. Physics-informed machine learning. Nature Reviews Physics, 2021, 3: 422-440.
\end{itemize}
\begin{itemize}
  \item  Kohtz, S., Xu, Y., Zheng, Z., \& Wang, P. Physics-informed machine learning model for battery state of health prognostics using partial charging segments. Mechanical Systems and Signal Processing, 2022, 172: 109002.
\end{itemize}
\begin{itemize}
  \item  Kendall, A., Gal, Y., \& Cipolla, R. Multi-task learning using uncertainty to weigh losses for scene geometry and semantics. Proceedings of the IEEE Conference on Computer Vision and Pattern Recognition, 2018: 7482-7491.
\end{itemize}
\begin{itemize}
  \item  Gal, Y., \& Ghahramani, Z. Dropout as a Bayesian approximation: Representing model uncertainty in deep learning. Proceedings of the 33rd International Conference on Machine Learning, 2016: 1050-1059.
\end{itemize}
\begin{itemize}
  \item  Lakshminarayanan, B., Pritzel, A., \& Blundell, C. Simple and scalable predictive uncertainty estimation using deep ensembles. Advances in Neural Information Processing Systems, 2017, 30.
\end{itemize}
\begin{itemize}
  \item  Zhao, H., Chen, Z., Shu, X., Shen, J., Lei, Z., \& Zhang, Y. State of health estimation for lithium-ion batteries based on hybrid attention and deep learning. Reliability Engineering \& System Safety, 2023, 232: 109066.
\end{itemize}

\end{document}
